\begin{textsample}{POS Dim 2 – ma – Score 31.00 – ma/ma\_em\_51.txt}  \label{ex:f2_pos_053}
5.
A educação \textbf{técnica} \textbf{profissional} e o 5º \textbf{itinerário} Ao longo da história, o ensino médio tem ocupado lugar de destaque na agenda das \textbf{políticas} públicas educacionais brasileiras em busca da melhoria na \textbf{qualidade} da educação oferecida aos estudantes nessa etapa da educação básica.
Com caráter a priori elitista, o ensino médio obteve democratização em sua \textbf{oferta}, porém sinalizada por uma dicotomia representada por diversas formas de escolarização \textbf{destinadas} a vários grupos sociais brasileiros ( RAMOS et al., 2011 ).
As últimas décadas expressam atenção e cuidado para a superação dessa dicotomia histórica, refletidos nas \textbf{diretrizes} e bases atuais da educação nacional estabelecidas pela Lei nº 13.415/2017, que altera a Lei de \textbf{Diretrizes} e Bases da Educação Nacional, a Base Nacional Comum Curricular, as \textbf{Diretrizes} Curriculares Nacionais para o Ensino Médio, \textbf{instituídas} pelo Conselho Nacional de Educação, e a \textbf{Portaria} nº 1.432/2018, que delibera sobre os Referenciais Curriculares para a Elaboração de \textbf{Itinerários} \textbf{Formativos}.
Como parte dessas reformulações, destaca -se a proposição de democratização da \textbf{oferta} da educação \textbf{profissional} e \textbf{técnica} integrada ao \textbf{currículo} do ensino médio enquanto uma das possibilidades de \textbf{itinerário} \textbf{formativo} que deve estar ao alcance das escolhas e interesses dos estudantes.
Contudo, os desafios para a materialização dessa proposta \textbf{formativa} são numerosos e exigem ações que envolvem desde o financiamento educacional até a capacidade física, organizacional e \textbf{técnica} do Estado para sua efetiva execução rumo à articulação entre trabalho, educação e práticas sociais, que assegurem tanto a construção e reconstrução de conhecimentos, habilidades, atitudes e valores essenciais às novas gerações quanto a concretização de objetivos pessoais, sociais, culturais, econômicos e \textbf{profissionais} ( BRASIL, 2018 ).
A \textbf{reforma} no ensino médio, objetivando a \textbf{garantia} do direito à educação mediante acesso e permanência dos estudantes na escola com \textbf{equidade} e \textbf{qualidade} educacional, prevê maior diversificação de percursos \textbf{formativos}, ao mesmo tempo em que reduz a parte \textbf{destinada} à formação geral básica do \textbf{currículo}, composta pelas quatro áreas do conhecimento, a saber : Linguagem e suas Tecnologias ; Matemática e suas Tecnologias ; Ciências da Natureza e suas Tecnologias ; e Ciências Humanas e Sociais Aplicadas.
Articulada à formação geral, a diversificação curricular está organizada por \textbf{itinerários} \textbf{formativos} com foco no aprofundamento das aprendizagens relacionadas às quatro áreas do conhecimento e à formação \textbf{técnica} e \textbf{profissional}, que corresponde à \textbf{oferta} do 5º \textbf{itinerário} \textbf{formativo}.
A proposta curricular do ensino médio do estado do Maranhão, no âmbito da \textbf{oferta} da educação \textbf{profissional} e \textbf{técnica} como um dos caminhos que o estudante pode trilhar durante sua \textbf{trajetória} opcional de formação nessa etapa da educação básica, está organizada a partir da integração dos quatro eixos \textbf{estruturantes} – investigação científica, processos criativos, mediação e intervenção sociocultural e \textbf{empreendedorismo} – estabelecidos no documento Base Nacional Comum Curricular ( BNCC, 2018 ).
Essa organização objetiva o desenvolvimento de competências preestabelecidas pela BNCC e habilidades necessárias à formação integral do estudante que, por sua vez, estão associadas a outras habilidades básicas requeridas pelo mundo do trabalho e pelas distintas ocupações, conforme estabelecem o \textbf{Catálogo} Nacional de \textbf{Cursos} \textbf{Técnicos} ( CNCT ) e a Classificação Brasileira de Ocupações ( CBO ).
A educação \textbf{profissional} e \textbf{técnica} integrada às modalidades e \textbf{ofertas} do ensino médio, no âmbito da \textbf{rede} de ensino público do Maranhão, compreenderá programas e \textbf{cursos} que promovam a formação e/ou \textbf{qualificação} \textbf{profissional} de estudantes para o desenvolvimento de \textbf{trajetórias} de vida e \textbf{carreira} \textbf{profissional} vinculadas às atuais ocupações, contextos econômicos locais e regionais e exigências da sociedade produtiva contemporânea.
Essas \textbf{ofertas}, considerando os eixos tecnológicos e orientações legais, poderão ocorrer compreendendo : - \textbf{Cursos} \textbf{técnicos} de nível médio integrados à formação geral básica ( em tempo parcial ou integral ) e/ou concomitantes ; e subsequentes, \textbf{destinados} aos egressos do ensino médio, previstos no \textbf{Catálogo} Nacional de \textbf{Cursos} \textbf{Técnicos} ; - \textbf{Cursos} \textbf{técnicos} em caráter experimental, com \textbf{cargas} \textbf{horárias} equivalentes a 800, 1.000 e 1.200 horas ; - \textbf{Cursos} de \textbf{qualificação} \textbf{profissional} : \textbf{cursos} de formação inicial e continuada ( FIC ) ; - Programas de aprendizagem \textbf{profissional} ; - Reconhecimento e \textbf{certificação} de saberes ; além de \textbf{certificação} por etapas com terminalidade.
Todas essas formas de \textbf{oferta} estão \textbf{amparadas} por \textbf{legislações} nacionais e \textbf{estaduais}, a saber : a Lei nº 13.415/2017, que altera as Leis nº 9.394/96 e nº 11.494/2017 ; a Resolução nº 3/2018, que \textbf{atualiza} as \textbf{Diretrizes} Curriculares Nacionais para o Ensino Médio ; a Lei de \textbf{Diretrizes} e Bases da Educação Nacional ( nº 9.394/96 ) ; a Resolução CNE/CP nº 1, de 5 de janeiro de 2021, que define as \textbf{Diretrizes} Curriculares Nacionais Gerais para a Educação Profissional e Tecnológica ; a Resolução nº 120/2013 ( CEE/MA ), que estabelece normas para a educação \textbf{profissional} \textbf{técnica} de nível médio no Sistema Estadual de Ensino do Maranhão ; e a \textbf{Portaria} nº 1.432/2018, que estabelece os referenciais para elaboração dos \textbf{itinerários} \textbf{formativos}, entre outras \textbf{regulamentações}.

% matched lemmas: amparar, atualizar, carga, carreira, catálogo, certificação, currículo, curso, destinar, diretriz, empreendedorismo, equidade, estadual, estruturante, formativo, garantia, horário, instituir, itinerário, legislação, oferta, política, portaria, profissional, qualidade, qualificação, rede, reforma, regulamentação, trajetória, técnico
\end{textsample}
